% ============================================================
\chapter{Design}
\label{chap:arch}
% ============================================================

\section{System Architecture}

PipedPeer is built around a long-lived \textbf{daemon} on every machine, which
handles discovery, admission, execution and result delivery; the
\texttt{pipedpeer} command is a short-lived client of it, and no state lives in
the command.

The property that distinguishes the system from the frameworks of
Chapter~\ref{chap:lit} is that \textbf{the coordinator is a library linked into
the client, not a service} (Figure~\ref{fig:architecture}). There is no scheduler process, no head node and no
master. A machine is submitter and worker at once, and the same binary does
both; each job is coordinated by the machine that submits it, and a training
ring's synchronisation by the daemon of its lead rank, chosen per job. A new
participant runs \texttt{pipedpeer connect} with the one-line code another
member's \texttt{pipedpeer invite} printed; on one network the machines also
find each other unaided. Across networks they meet through an
\textbf{introducer}, an address book that tells machines where to find each
other and carries no job data. Every daemon is an introducer for its own
cluster. Two optional roles exist: \texttt{chunks-only}, for a phone that takes
parallel chunks but not whole jobs, and \texttt{introducer-only}, for a hosted
server that runs nothing.

\begin{figure}[!htbp]
    \centering
    \includegraphics[width=150mm]{\figpath{fig_3.1_system_design.png}}
    \caption[PipedPeer system design]{PipedPeer system design. The submitting side builds the environment
    and places the job; the worker admits it, caches the closure and runs it
    sandboxed under a memory ceiling. Both roles are the same binary. Peers on
    other networks are reached over direct QUIC links that an introducer helps
    establish but never carries; every daemon is an introducer for its own
    cluster. Phones take parallel chunks.}
    \label{fig:architecture}
\end{figure}

The principal subsystems are listed in Table~\ref{tab:subsystems}. All are
packages under \texttt{src/internal/}; the remainder are supporting code
(process management, setup, paths, logging).

\begin{table}[!htbp]
\centering
\caption{Subsystems and their responsibilities}
\label{tab:subsystems}
\renewcommand{\arraystretch}{1.2}
\rowcolors{2}{}{rowgray}
\small
\begin{tabularx}{\textwidth}{L{3.2cm} X}
\toprule
\textbf{Package} & \textbf{Responsibility} \\
\midrule
\texttt{coordinator} & Filters, scores and selects a node; takes the lease; retries on failure \\
\texttt{daemonapi} & HTTP and WebSocket surface, lease manager, sandbox, pool executor, gradient reduction, closure caches, node poller and address gossip \\
\texttt{schedule} & Makespan-optimal division of a chunk across measured devices \\
\texttt{ddpplace} & Choosing the ranks of a training ring by measured throughput \\
\texttt{discovery} & UDP multicast and broadcast probes, TCP health sweep fallback \\
\texttt{internet} & Peer manager across networks: register with every introducer, punch, link, migrate, forward; the in-daemon introducer \\
\texttt{direct} & Shared UDP socket, candidate race, birthday-collision fallback, QUIC links \\
\texttt{rendezvous} & The introducer's address book and punch-request forwarding \\
\texttt{portmap} & Router port mapping by PCP, NAT-PMP and UPnP-IGD \\
\texttt{authtoken}, \texttt{tlsid} & Shared-secret middleware; certificate from the node key, pins by node id, TLS and HTTP on a single port \\
\texttt{clustercfg}, \texttt{invite} & The one configuration file (secret, introducers, role); the invite code \\
\texttt{identity}, \texttt{nixsign} & Ed25519 node key, whose fingerprint is the node id; signing closures and the trusted-key list \\
\texttt{narpack}, \texttt{nixstore} & Signed closure archives and cache pruning; system or private (rootless) Nix store \\
\texttt{cgroups} & Memory ceilings for the daemon and for each job \\
\texttt{nodestore} & SQLite address book: every path to every node, last observed capacity \\
\texttt{nixgen}, \texttt{pythondeps} & Flake generation, wheel pins and the interception shim; import scanning and resolution \\
\texttt{resourceest}, \texttt{runrec} & Tiered memory estimate; the run record of timed steps \\
\texttt{heartbeat}, \texttt{gpu} & Capability and load reporting; vendor GPU detection \\
\texttt{service}, \texttt{warmup}, \texttt{doctor}, \texttt{phone} & Autostart unit; background environment builds; diagnostics; phone worker \\
\bottomrule
\end{tabularx}
\end{table}

\section{Data Flow Diagram}

\subsection*{Level 0 --- Context Diagram}
\addcontentsline{toc}{subsection}{Level 0 --- Context Diagram}

At the highest level (Figure~\ref{fig:dfd0}) the system has three external entities: the
\textbf{user}, who submits a script and receives output and result files;
\textbf{peer nodes}, which supply capacity and return computed results; and an
optional \textbf{introducer}, which exchanges candidate addresses so that peers
on other networks can reach one another.

\begin{figure}[!htbp]
    \centering
    \includegraphics[width=150mm]{\figpath{fig_3.2_dfd_level0.png}}
    \caption[Level 0 DFD (context diagram)]{Level 0 DFD --- context diagram. The introducer is dashed because
    the system works without it on one network.}
    \label{fig:dfd0}
\end{figure}

\subsection*{Level 1 --- System Decomposition}
\addcontentsline{toc}{subsection}{Level 1 --- System Decomposition}

\begin{figure}[!htbp]
    \centering
    \includegraphics[width=146mm]{\figpath{fig_3.3_dfd_level1.png}}
    \caption[Level 1 DFD]{Level 1 DFD --- seven processes over four data stores.}
    \label{fig:dfd1}
\end{figure}

Figure~\ref{fig:dfd1} decomposes the system into seven processes:

\begin{itemize}[leftmargin=*]
  \item \textbf{1 Resolve environment.} Scan the script's imports, generate a
  flake, build the closure.
  \item \textbf{2 Find and link peers.} Probe the subnet; across networks,
  publish candidates to the introducer, punch, and register each direct link as a
  node.
  \item \textbf{3 Place task.} Apply hard filters, score the survivors, reserve
  a lease on the best.
  \item \textbf{4 Transfer.} Upload the workspace, and of the closure only the
  store paths the worker lacks; signed between daemons.
  \item \textbf{5 Execute.} Run inside a \texttt{crun} sandbox under a memory
  ceiling, streaming output back as it is written.
  \item \textbf{6 Distribute parallel work.} Share chunks among peers by
  measured throughput, subject to the cost model.
  \item \textbf{7 Return results.} Send back changed files and a deletion
  manifest.
\end{itemize}

The four stores are the address book, the closure cache, job history, and the
identity and trust store (node key, pinned certificates, trusted Nix keys).

\section{Use Case Diagram}

Figure~\ref{fig:usecase} shows who uses the system and for what; machines are
actors as well as people.

\begin{itemize}[leftmargin=*]
  \item \textbf{User:} runs a script on the cluster, asks for a GPU or a memory
  floor, trains with data parallelism, and inspects jobs and nodes.
  \item \textbf{Node operator:} sets up a machine, invites another machine, and
  connects a machine with an invite code.
  \item \textbf{Peer daemon:} accepts and leases work, serves signed closures,
  executes parallel chunks, and introduces its cluster's peers to one another.
  \item \textbf{Phone (Termux):} joins under \texttt{proot} and Nix and executes
  parallel chunks.
  \item \textbf{Hosted introducer:} registers nodes and forwards punch requests;
  it carries no job data.
\end{itemize}

\begin{figure}[!htbp]
    \centering
    \includegraphics[width=129mm]{\figpath{fig_3.4_use_case.png}}
    \caption[Use case diagram]{Use cases. A peer daemon is an actor in its own right, because
    every node serves closures, executes chunks and introduces its cluster's
    peers; a phone executes chunks; a hosted introducer only registers nodes
    and forwards punch requests.}
    \label{fig:usecase}
\end{figure}

\section{Sequence Diagram}

\begin{figure}[!htbp]
    \centering
    \includegraphics[width=120mm]{\figpath{fig_3.5_sequence.png}}
    \caption[Sequence diagram of one job]{End-to-end execution of one job.}
    \label{fig:sequence}
\end{figure}

Figure~\ref{fig:sequence} follows one job from submission to results:

\begin{enumerate}[leftmargin=*]
  \item The user runs \texttt{pipedpeer run script.py}.
  \item The client scans imports, builds a Nix closure locally, and estimates
  the job's memory requirement.
  \item It asks its own daemon for the healthy peers, including peers on other
  networks, reached through local forwarders onto direct links.
  \item The coordinator filters and scores them and sends \texttt{accept} to the
  best candidate, which checks its capacity and reserves it.
  \item The client sends \texttt{commit}, which turns the reservation into a
  running lease.
  \item It asks whether the worker can already run the closure; if not, it asks
  which store paths are missing and sends only those, then uploads the
  workspace.
  \item The worker pushes the closure to peers that lack it, so that they can
  take a share of the parallel work.
  \item Execution proceeds inside a sandbox under a per-job memory ceiling, with
  output streamed back as it is written.
  \item If the script uses a parallel primitive, the shim distributes it subject
  to the cost model.
  \item On exit, changed files and a deletion manifest return as a compressed
  archive and are extracted into the user's directory.
\end{enumerate}

Every request in this sequence carries the cluster's shared secret and travels
over pinned TLS on a local network or inside a QUIC link across networks.
